\section{结果}

\subsection{具有明确来源信息的文献关联人类 TE 图谱}

当前 Neo4j 快照包含 225 个 TE 节点和 2,308 个 Paper 节点，同时包含 12,444 条连接 TE 与生物医学实体的有向 \texttt{BIO\_RELATION} 关系。当前每条生物学关系均具有标准化谓词、至少一个 PMID 字段和支持 PMID 计数。这些属性使界面中的边可以追溯到用于支持该关系的文献记录，但可追溯性并不会把提取出的关联转化为因果证据。

生物医学模式包含 11 类实体。目前数量较多的标签包括 Function（3,683 个节点）、Gene（1,280）、Protein（1,089）、Disease（676）、RNA（588）、Mutation（377）和 Pharmaceutical（293）。数量较少的类别包括 Toxin（67）、Lipid（26）、Peptide（23）和 Carbohydrate（12）。DiseaseCategory（767）、Paper（2,308）及分类相关节点单独表示，因为它们承担组织或来源追踪作用，而不属于上述 11 类生物医学实体。

\begin{table}[htbp]
\centering
\small
\caption{当前 TE-KG 内容快照。各项计数代表不同的运行时单位，不得相加。快照于 2026 年 7 月 31 日完成验证。\label{tab:snapshot-zh}}
\begin{tabularx}{\textwidth}{@{}p{0.20\textwidth} r Y Y@{}}
\toprule
组成部分及单位 & 数量 & 定义 & 来源与边界 \\
\midrule
Neo4j TE 节点 & 225 & \texttt{tekg3} 图中的 TE 实体 & 实时标签计数查询；与 Browse 条目不是同一单位 \\
Neo4j Paper 节点 & 2,308 & 图中表示的保留论文 & 实时标签计数查询；与最终文献语料库一致 \\
有向生物学关系 & 12,444 & 已存储的 \texttt{BIO\_RELATION} 关系 & 按存储方向计数一次；当前全部记录均含谓词和 PMID 来源字段 \\
生物医学实体类别 & 11 & 碳水化合物、疾病、功能、基因、脂质、突变、肽、药物、蛋白质、RNA 和毒素 & 不包括 Paper、DiseaseCategory 和分类标签 \\
Browse 目录条目 & 276 & 版本化 TE 目录记录 & 由 MySQL 支持的 Browse API；与图 TE 节点属于不同记录单位 \\
已分类 TE & 225 个中的 192 个 & 实时摘要中已具有分类类别的 TE 节点 & 由 Neo4j 支持的分类 API \\
表达样本 & 1,158 & 205 个正常组织、307 个正常原代细胞和 646 个癌细胞系样本 & 三类来源组；不是跨情境配对队列 \\
可检索共表达条目 & 784 & 285 个 TE 和 499 个基因条目 & 三类情境的经批准展示目录；不是完整离线网络的特征数量 \\
下载文件 & 10 & 6 个表达文件、2 个图文件和 2 个分类文件 & 当前网页入口；仍需公开归档发布 \\
\bottomrule
\end{tabularx}
\end{table}


\subsection{分类视图与目录视图保留不同的记录单位}

实时分类摘要包含 225 个 TE，其中 192 个已分配分类类别。摘要还报告 188 个标准叶节点，并将包含 36 条记录的全来源树与包含 189 条记录的 RMSK 加 RepBase 树分开。这些数值描述的是 API 定义的分类视图，而不是彼此独立的生物学总数。Browse 目录包含 276 条版本化记录。由于 Browse 与 Neo4j TE 节点代表不同的记录单位，该较大的 Browse 数量单独报告，不能与图中的 TE 节点数合并为一个 TE 总数。

分类界面在使用同一套 API 分类数据的前提下，同时提供层级树和力导向图。图例可以暂时隐藏或恢复 TE 类别，而不会改变底层结果集。这种交互有助于查看密集分类，但不会引入第二个分类事实来源，也不会修改分类归属。

\subsection{表达与共表达增加三类生物学情境}

当前表达矩阵包含 205 个正常组织样本、307 个正常原代细胞样本和 646 个癌细胞系样本。共计 1,158 个样本，为 TE 和基因表达查看提供三类情境。由于这些情境来自不同研究和采样设计，系统将其相互分开。当前数据不支持把它们视为配对的实验比较。

正式共表达目录覆盖相同的三类情境，并包含 285 个 TE 和 499 个基因的经批准可检索条目。三个情境合计提供 849 条展示建议：17 条标记为核心案例，287 条标记为高置信条目，542 条标记为可检索条目，另有 3 条不建议默认展示。这些层级用于控制呈现和可检索性，并不代表统计显著性类别。所有展示边均已满足固定的相关系数和 FDR 阈值，单次返回的展示网络只是更大离线结果的受限视图。

\subsection{相互连接的界面提供互补证据入口}

TE-KG 针对不同问题提供多个访问入口。Browse 整合了 TE 目录筛选、内置搜索以及查看所选 TE 结构化记录的功能；Knowledge Graph 展示文献关联的生物医学关系并提供 PubMed 证据；Path 搜索所选实体之间的连接；Expression 展示特定情境下的丰度谱；Co-expression 展示 TE-基因相关邻域；Download 当前提供 10 个文件，包括 6 个表达文件、2 个图文件和 2 个分类文件。

这些入口在 TE 身份层面相互连接，但保留不同的证据含义。例如，打开图中的疾病关联边可以访问支持论文，而不是声称该 TE 导致疾病。同样，打开共表达邻域会报告分析情境和相关阈值，而不是调控机制。保持这种区分是面向用户的设计组成部分，而不仅是内部实现细节。

\subsection{自然语言追问仅限当前会话}

Agent 和 DeepThink 可以根据英文问题检索多个本地证据层，并理解当前对话中的追问。生成回答中的 PubMed 标识符会显示为指向 PubMed 的链接。写作层使用普通科学语言表达证据限制，不向用户展示内部路由、置信标签或其他开发诊断信息。例如，回答可以将一条关系描述为仍需实验验证的关联，而不要求读者理解内部状态代码。

当前评估支持功能覆盖结论，但不支持总体准确率结论。固定的 13 题基线覆盖了全部 12 类检索角色，后续回归检查在定向修复后保护了图路由、文献链接和多插件行为。更早的一次 36 题测试包含后来已经修复的失败，但整个测试集尚未作为同一版本锁定的基准重新运行。因此，本文将该界面描述为受边界约束的检索入口，并保留具体已知限制，而不把历史通过率作为当前性能发表。

\subsection{跨证据层的 L1HS 使用案例}

L1HS 被选为主要示例，因为它可以在当前分类、图、序列、基因组、表达和共表达界面中解析。最终图件将在每个分面中保留结果类型的区别，并将所有文献陈述连接到已核实的 PMID。\draftnote{本小节达到论文可用状态前，需冻结准确的 L1HS 记录、表达来源表、共表达邻域，以及一到两条由 PMID 支持的图关系。不得使用通用模型知识填补缺失证据层。}
